\section{Scaligero Meets Evola}

\subsection*{First part}
\begin{quotex}
In this essay, \textbf{Massimo Scaligero} describes his first impressions on meeting \textbf{Julius Evola}. Although they went in different directions (Scaligero became an Anthroposophist), they remained good friends and collaborated for several years. Like Evola, Scaligero was interested in the Roman spirit in the context of spiritual races. Scaligero incisively notes the essential differences in their worldviews. Scaligero uses thinking and imagination to penetrate into reality, and ultimately to the Logos. Evola, on the contrary, insists on the primacy of will.

\end{quotex}

I knew Evola at a time when almost everyone prudently distanced himself from him: he revealed himself, by means of the journal \emph{La Torre}, to be the boldest protester of the culture of the Regime. Even though those around him created a type of vacuum, for that reason he carried on the attack unperturbed. That attracted my attention, even if I did not yet understand what he was really driving at. In those days, around the spring of 1930, I went to him despite being advised against it by well-meaning friends: I did not want others to decide for me. At that time telephones in the home were rare, so the appearance at someone's home would not have been considered an indiscretion: in case of the absence of the person you wanted to visit, we used to leave a calling card with the concierge.

I was curious to know the person who showed so much courage and about whom Adriano Tilgher had spoken to me with admiration, albeit with reservations about the content of his doctrines. I remember he called him ``the most powerful dialectician of Europe''. But that was not what impressed me. I had yet to understand that Tilgher was referring mainly to the dialectical battles at the \emph{Associazione per il Progresso Morale e Religioso} (Association for Moral and Religious Progress), then directed by professor Giuseppe Puglisi, and headquartered in Piazza Nicosia (later, it was transferred to Piazza Campicelli under the direction the poet Raniero Nicolai) where, at the end of the conference, a debate took place in which the usual participants included Professor John Imperato, Professor Ernesto Quadrelli, Professor Giuseppe Palanga, sometimes Tilgher, and, finally, Evola, who regularly beat them all with compelling arguments and humorous barbs.

I knocked on the door of the next to top floor of Corso Vittorio 197 and a youthful, tall, rangy character, undoubtedly older than I, opened the door: his expression was between Buddhist and Olympic, his bearing quite calm. Having at once intuited the meaning of my visit, or rather, lack of one, Evola promptly met me with genuine affection and this affection was the strength of the connection, beyond dialectics and doctrine, that linked me to him for years.

As I wrote in my book \emph{From Yoga to the Rosicrucians}, I didn't know that Evola had followed an esoteric path for a time and directed a group of aspirants to initiation, although his name was familiar to me from the journal \emph{Ignis}.

I was already following my own very personal path since adolescence, which was between Yoga and the most unrestrained thought of the West --- Nietzsche, Stirner, Bauer --- but I did not like to talk about it with anyone.

I recall that at first Evola was surprised by the meaning of my visit, neither political, nor esoteric: as I said, at that time he was subject to continuous attacks and intimidations. I remember that the theme of mountains and the inner impressions of the climb, the silence and solitude of the peaks, united us right away. In a chapter devoted to Evola, in the book I mentioned, I outline the meaning of my relationship with him: the encounter with a world that gave me the perception of thought as a primordial force.

I loved Goethe and Nietzsche for that reason: not the Nietzsche of the exaltation phase, but rather that of \emph{The Birth of Tragedy} and \emph{Thus Spoke Zarathustra}.


At that time, I dramatically felt the state of death of the thinking of the whole culture, every expression of culture, whatever its subject, or its political complexion: I had the impression of moving in a vast cemetery: I came out of it every time thanks to energetic meditation, but it was important for me to encounter, through Evola, a kind of thinking that I secretly cultivated in myself: a thinking still capable of life and freedom. After several years, moreover, I had to believe that it was not so much the content of Evola's writings which delighted my inner self, but also feeling that content as a very personal construction: a work of art: a work of brilliant, organic art, even when I felt the affects --- to be clear --- of Tantrism of the ``left Hand''.

The harmony with Evola was essentially this: the contact with a world of forces. But this was possible because I myself carried in me this world of forces: the recognition of their origin, however, gradually brought about the orientation that was to carry me out of the path. But what is not in dispute: reality is one, truths are many.

The dialogue with Evola reinforced in me the need for a distinction of values and traditions: that was decisive for me with regard to issues such as karma and reincarnation, the essentially Christian character of Western alchemy, the mysteriously Christian character of the Grail, the fundamentality of the Logos as Christ-Principle. I understood Evola's paganism, but I considered it only as a positive impulse of the quest for what really was the Logos beyond all worldly expressions and beyond what He had become in theological and ethical-political formalism.

Actually Evola is a strong medicine, a shock treatment, the application of force without mediation: as such, he acts in those who already have the power of mediation and just has the task of making it operative. But if he truly realized that, he already has the principle of the Force. Thanks to Evola, I encountered a decisive moment of my inner experience: the feeling of liberation, which gives the image rather than knowledge of liberation: the heroic, mythic image, not cognitive. I then had to discover that the cognitive way was necessary for me: Evola's way had the cognitive appearance, but in essence it relied on the feeling of strength, not on its reality, penetrable and comprehensible by the power of the idea. To renounce this power of the idea, means to accept super-nature mystically, or by faith.

In other words, we do not enter into the magical by the feeling of magic or by the dialectic of magic, but only by virtue of its initial power in the soul, which is ideational or imaginative power. Evola, indeed, cannot hide a certain contempt for thinking, or for a ``way of thinking'', believing thinking as the expression of a specific sentient nature and not of super-nature.

The ideal of the absolute individual was for him an expression of the will sufficient in oneself: thinking indeed has nothing to do with it, and he is right, but you cannot move away from it: nor can you willfully experiment with thinking, which ceases to be dialectical, but volitional. He pointed out that in yoga, essentially the strength of the will was awakened and recovered in its magical movement, and it becomes the flow of Kundalini.

The whole power of evolian ``persuasion'' is this immediacy of the possessed will. He refuses, as premise, thinking, the cognitive act, the mediation of the idea: Evola always speaks about an absolute act, a movement from which they come about, as it should be in itself magical.

It is in effect the raw material of the Work, which it is necessary to already have: Evola attributes it to a particular inner constitution, to belonging to a particular race of the Spirit. For me it was always clear that it is about the meeting of the free I with its own karma: This is the condition of the element with which a psychophysiological I covers itself.

Nevertheless, that persuasion, so decisively expressed by Evola, was a confirmation of my personal experience of the centrality and the invincibility of the I: which, brought to the depths, can encounter the Logos, His origin: Evola was not able to recognize that. I knew that the formulation of the centrality and asceticism of the I was correct: I appreciated Evola, because at that time I was convinced that he had never come to talk about it so explicitly and methodologically. Stirner's and Nietzsche's attempts were ultimately Hegelian-romantic: bold offshoots of the Hegelian left, but prisoners of dialectics.

\subsection*{Second part}

\begin{quotex}
This is the conclusion of \textbf{Massimo Scaligero}'s essay describing his first impressions on meeting \textbf{Julius Evola}. In this part, Scaligero gets to the heart of Evola's thinking, describing what unites them and what separates them.

Unlike other reviews that describe Evola as a political figure, a philosopher, etc., Scaligero understands Evola in the way that one esoterist can understand another.

\end{quotex}

The way which I personally followed at that time was my own synthesis, which, I believed, could not correspond to a doctrine. The term ``Science of the I'' resonated pleasantly familiar with me: in fact it corresponded to the ideal of asceticism to which I tended and to the correlative method. This was the true help that I got from Evola: to perceive the doctrinally justifiable asceticism of the I, independently of Tradition. The reality of this proved to be traditional, even as it did not correspond to any conceivable form in Tradition. In a similar way, however, it was possible to come to see the source of this asceticism. For a while I had the feeling of having found in Evola the master that I was looking for, and he had indeed more than a claim.

I think I was affected more than Evola, by the disappointment of the fact that, at a certain time, even as I was following the ``way of the I'', I was no longer able to follow his path. I could recognize in him only a characteristically strong indicator, an individuality from autonomy so sufficient in himself, since it does not require \emph{metanoia}, or conversion, of his own dialectic to arouse in others the feeling of his strength through dialectic: naturally, without possibility of the resolution of individual limits: in the end a resolute titan, whose strength however had its foundation in himself, but he could not become the weakness of the less autonomous disciple who is unable to find the foundation in himself, that is, his inner path.

Along the way, I could understand my relationship with Evola even better. Generally it happened --- and it is happening --- that one reads his works and then decides to follow him. In substance it had happened to me in a something different way: I had gone to him without reading a single line of his: for a long time our conversations had as its theme the mountains, nonconformity, the courage to oppose the pseudo-culture. I kept company with a group of aggressive transtiberine friends, that I mobilized when he was in danger of being attacked again on the occasion of the different causes that he had in court for lawsuits and the aftermath of lawsuits. I remember that, at our second meeting, he lent me Meyrink's \emph{Golem}, probably to give me the opportunity for an esoteric connection with him; but, just flipping through the pages, I was not too convinced of their content. In truth I did not read the book at all and after a fortnight gave it back to Evola, avoiding telling him that I was not interested. I noticed, however, in Evola a minimal expression of disappointment. Later, however, I carefully read Evola and Meyrink.

The true meaning of what had driven me inwardly toward him, became clear to me with the passing of time: it was the meeting with expert disciples of Steiner, collaborators of Ur (although I was already connected with Giulio Parise and Arturo Reghini through the journal \emph{Ignis}), among whom the significant and decisive one for me had to be Colazza, and the inner encounter with the one whom I came to recognize as the Master of the new age, Rudolf Steiner; in fact, the original edition of Ur contained impartial appreciations of Steiner. Actually, it was not so much the knowledge of Evola's works that united me to him for several years, but rather his lively, exceptionally individual, personality. Of course, I then would examine these works line by line, before making the decision for another path.

The meetings I had with Evola were filled with a Hermetic climate, from which the feeling of magical forces regularly arose, something like what I actually had from the habit from personal meditation. Or else there were Bacchic-gastronomic meetings, based on spirited philosophy and bucatini all'Amatriciana\footnote{\url{https://en.wikipedia.org/wiki/Amatriciana_sauce}}, but with the unusual addition of peperoncino. I was struck by the fact that Evola, who was a compelling alcohol supporter, cared to drink very sweet wines such as vino passito\footnote{\url{http://italian-flavor.com/notes/passito.php}}, or Gaglioppo\footnote{\url{http://drinks.seriouseats.com/2011/01/calabrias-gaglioppo-grape-wine-southern-italy.html}}. He explained it to me, saying, I don't know how jokingly, that, having learned from Colazza that alcohol deadened the I, while sugar was the physical food of the I, he, not to take risks, tried to reconcile them. Moreover, he connected me to him by an inclination that has always guided my action, to support the one against whom everyone, the \emph{koinonia ton kakon} communion of evil , manages to feel themselves sympathetic: so-called ``public opinion'', or Phariseeism, whose hypocrisy I have always felt. At that time Evola had only adversaries, even if, truth be told, he did everything to provoke them.

The idea that the truth was from the party against whom everyone lined up in solidarity, in a certain sense corresponded for me to a specific supersensible law. On the other hand, Evola today is translated into the major languages of the world, and he has a following even in Asia. However Evola never assumed the role of the ``master'': guidance counselor, instigator, yes, but not the master. Sometimes when I had problems with respect to the method, and I turned to him, he demurred, recommending Colazza or Bonabitacola to me, that is, Steiner's disciple or Kremmerz' disciple. Evola found Bonabitacola interesting especially because he received all insults like a Zen master; while Colazza was the character for whom he professed the utmost deference, because he recognized him as an outstanding investigator of the ``occult''.

Colazza also had respect for Evola as a strong personality: he reproached him only for his yoga and the fact that he wrote books. ``One of us'', he said to me one day, ``must not write books. Books occultly bind the author, they bind him to his present thought, preventing him from opening himself up to the new, to the unknown, to what he still is not capable of thinking. Besides, whoever plays the role of teacher, loses the opportunity to be a disciple.''

Something substantially true can be recognized with this thought of Colazza. Nevertheless, you cannot help but consider in Evola's case the mission of sacrifice by the word: namely the compromise of one's own way, on the dialectical level, for the way of others. The snare that he throws would not be necessary if the initiatory path were not in relation to a free choice, to self-determination, that is a determination that is not fated to occur, otherwise it would not be free. This can sufficiently explain his dominant and bold reconciliation of Buddha with Nietzsche; for us it is valid, even if only introductory to the true sadhana. When esoterism is realized, it ceases to be mental position.

The enigma of Evola'a inner personality and his relationship with Tradition remains. In truth, the traditional form cannot conceal the powerful thrust of his anti-traditional system of thought. If you observe, Evola makes use of the traditional element for building his own absolutely personal spiritual cosmos. He uses the value of Tradition as the aristocratic touchstone, that is, as the opposite to the modern world. He seems to propose a return to Tradition, but in fact, he wants something that is not Tradition, indeed positively against it. And this is what we can recognize as important in Evola.

Guenon is in Tradition, Evola leaves it all the time, while appealing to it: it is a juxtaposition, a dialectical theme, a \emph{habitus cogitandi} habit of thinking. From Novalisian ``magical idealism'' to Tantrism, from the Nietzschean formulation of Mahayana to the ``pagan'' interpretation of Alchemy and the Grail, to the sympathy for Michelstaedter and Meyrink, Nietzsche, Kremmerz and Crowley, the imperious authority of thinking is clear, that makes everything subject to an intimate personal vision, to an individual, determined ``will to power'', ultimately, to an impulse of freedom that subordinates everything to self-affirmation.

Evola succeeds in making his own world real through the power of thought. This is the aspect that could give real help to those who want to follow his path, not as a follower dazzled by his light, but as a conqueror of the Light. It does not matter whether the thinking is correct or not, it becomes correct when it is powerful. It is undeniably true: error, immorality, is weak thinking, or merely dialectical, or reflexive. The power of thought is the truth: the truth that coincides with reality. But only a free man can realize it and make it his own. Those who believe in that thought as a content that comforts him, excites him or hypnotizes him, is not free, he is not in the current of reality. He believes he is in Tradition, but he is against it in the worst form, that of dogmatism and mystical subjection. So we said, at the beginning of this essay, that Evola is a strong medicine, a shock treatment: a discipline that it is necessary to possess. If you do not possess yourself, you actually betray yourself, because you act with respect to it as one who is subjugated.

The living element of Evola's work is the idea of liberation, which requires, however, pure recognition, the decision of self-knowledge, by the person who encounters it. This is the demand and the virtue of his system, but also the real thrust of his personality, which can be seen as a synthesis of what he is individually with that which the real Guides of humanity implant in him, despite himself. Certain inspirations in the ascetic become operatively individual, even as they draw on the Super-individual. Julius Evola points out a direction that, to become creative in the esoteric sense, requires its separation from its phenomenology, that is, from \emph{maya}, from the \emph{ethos} resulting from it in a social, political sense: especially from this. Certain mixtures betray the spiritual proposition: the worst disasters always come from the collusion of the Sacred with the profane. A similar separation, just like the assumption of the discrimination of
\emph{subtile a spisso} the subtle from the dense, is about Evola's own karmic direction. You can metaphysically become aware of his cosmic-spiritual future. It will be determined by how much independence he has achieved from the written works. The possibility will be decisive that he not be identified with his own doctrinaire expression, i.e., with the mythical and mystical world of his followers.



\flrightit{Posted on 2015-03-10 by Aeneas}

\begin{center}* * *\end{center}

\begin{footnotesize}
\begin{sffamily}
\texttt{White Rabbit on 2015-03-11 at 23:56 said:}

It is interesting how Evola's path of negation and tearing down ideas can also be described as one of faith. For a thinker I'm sure it would require a strong intellectual propulsion from the start, at least.

Scaligero says that that way wasn't for him. A puzzle: is someone who tends to ``question everything'' fitted for the cognitive way, or Evola's way of demolishing all intellectual supports?

\hfill

\texttt{David R on 2015-03-12 at 22:21 said:}

Thank you for this translation. There is alot to pounder about. Might I ask, if you know, if Scaligero is of any worth in himself ? Does this text renders him justice ? Thank you.

\hfill

\texttt{Tom Blanchard on 2015-03-13 at 11:20 said:}

Scaligero has a truly unique perspective on Evola. His idea of Evola as a kind of ``titan'' seems accurate -- Evola frequently mentioned that Olympian heroes often came from titanic stocks/backgrounds, almost as titanism was a sort of prior phase or necessary pre-requisite state that nevertheless had to be wholly overcome and transformed.

Scaligero's description of Evola as a man whose style and approach was a ``necessary discipline'' but whom he nevertheless had to leave behind (without repudiation) is slightly reminiscent of Dante's relationship with Virgil in the Divine Comedy. This is even reinforced by his reflection that he for some time thought of Evola as an appropriate ``master''.

\hfill

\texttt{Cologero on 2015-03-13 at 23:41 said:}

David R: As Scaligero wrote, for a time he followed Evola, but then left him to develop his own idiosyncratic brand of Anthroposophy. We'll be making use of some of those earlier writings over time. As for his second phase, that is up to your own interests; I know that the Anthroposophists think highly of him.

\hfill

\texttt{Aeneas on 2015-03-13 at 23:47 said:}

Mr Blanshard, not only was Evola a ``master'' but according to Scaligero he also had spirit Guides. The Comedy is a good analogy; Evola just couldn't ``get'' the Logos.

Am I the only one who was more interested in the pasta and wine? BTW, the pasta recipe should be made with bucatini rather than spaghetti. If anyone makes it to South Beach, look me up and perhaps I'll cook it up.

\hfill


\end{sffamily}
\end{footnotesize}

